\documentclass[10pt,a4paper]{amsart}
\usepackage[margin=19mm]{geometry}
\usepackage{amsmath,amssymb,mathtools}
\usepackage[scr=rsfs,cal=euler]{mathalfa}
\usepackage{xfrac}
\usepackage[unicode,hidelinks,bookmarks=false]{hyperref}
\newcommand{\ie}{i.e.\ }
\newcommand{\cf}{cf.\ }
\newcommand{\resp}{respectively\ }
\newcommand{\Subsection}{\subsection}
\providecommand{\nobreakdash}{\nobreak\hbox{-}}
\setlength{\emergencystretch}{2em}
\multlinegap0pt
\usepackage{bm}
\usepackage{bbm}
\usepackage{dsfont}
\usepackage{xspace}
\usepackage{cancel}
\usepackage{mathtools}
\usepackage{amsthm}
\makeatletter
\@ifundefined{thm}{\newtheorem{thm}{Thm}}{}
\@ifundefined{lemma}{\newtheorem{lemma}[thm]{Lemma}}{}
\makeatother
\newcommand\jdz[1]{ \left| #1 \right| }
\newcommand\lk[1]{ \left\{ #1 \right\} }
\newcommand\norm[1]{ \left\| #1 \right\| }
\newcommand\sk[1]{ \left( #1 \right) }
\newcommand{\K}{{\mathcal K}}
\newcommand{\Q}{{\mathcal Q}}
\newcommand{\e}{\varepsilon}
\newcommand{\real}{{\mathbb R}}
\renewcommand\d{{\rm d}}
\title[Best constants in the vector-valued Littlewood-Paley-Stein t]{Best constants in the vector-valued Littlewood-Paley-Stein theory}
\author{Guixiang Hong, Zhendong Xu, Hao Zhang}
\date{Source: arXiv:2401.13932v1 (CC BY 4.0)}
\begin{document}
\maketitle

\noindent\emph{Source and licence.} Best constants in the vector-valued Littlewood-Paley-Stein theory. Version arXiv:2401.13932v1 (CC BY 4.0), distributed under the Creative Commons Attribution 4.0 International licence, as recorded in the arXiv licence metadata for this exact version and shown on the abstract page https://arxiv.org/abs/2401.13932v1. Full rights notice: licenses/arxiv-2401.13932-CC-BY-4.0.txt.\par\smallskip

\noindent\textit{Editorial scope.} Counted unit: the Lemma whose source label is \texttt{l3.8} in the article (source0.tex lines 1364--1373, source label \texttt{l3.8}), delivered together with the article's own complete original proof of it (source0.tex lines 1375--1408, one \texttt{proof} environment of 34 lines). No mathematics is written, repaired or re-derived: statement and proof are the article's own text line for line. Required context retained uncounted: derpt (lines 479--481); BoundTent (lines 1096--1109). Every definition, display and label of those lines is retained unchanged. Omitted from this package and not redistributed: 1--478, 482--1095, 1110--1363, 1374--1374, 1409--2575. Nothing inside a retained excerpt is omitted or altered: every retained body is delivered exactly as the article's lines are written, so no omission receipt of any kind accompanies this package. Citations: the retained excerpts carry no citation command at all, so no bibliography entry of the article is transcribed and no reference is redistributed. Reference adaptation: none was needed and none was made; every reference inside the delivered unit resolves inside it, so no unresolved reference or citation is produced. Printed slips kept literal: no printed word, symbol, hypothesis or number of the counted statement or of its proof was altered. Layout and class: the article's own class is \texttt{amsart} with packages including inputenc, geometry, mathabx, esvect, babel, mathtools, amstext, amsthm, amssymb, enumerate, esint, graphicx, bbm, xy; the wrapper is the lane's pinned \texttt{amsart} editorial wrapper at 10pt on A4, which restates the theorem-like environments the retained text needs and adds the article's own macro definitions that the retained text needs (\texttt{\textbackslash K}, \texttt{\textbackslash Q}, \texttt{\textbackslash d}, \texttt{\textbackslash e}, \texttt{\textbackslash jdz}, \texttt{\textbackslash lk}, \texttt{\textbackslash norm}, \texttt{\textbackslash real}, \texttt{\textbackslash sk}), with extra wrapper packages (bm, bbm, dsfont, xspace, cancel, mathtools). The delivered numbering is the unit's own. Assets: no retained span contains a figure environment, a graphics-inclusion command, a PDF-page inclusion, a table or a TikZ picture; no image file is copied into this package, the package declares no asset dependency and no image file is redistributed with it. Licence: version arXiv:2401.13932v1 of this article is distributed under the Creative Commons Attribution 4.0 International licence, as recorded in the arXiv licence metadata for this exact version and shown on the abstract page https://arxiv.org/abs/2401.13932v1. A full rights notice is delivered at licenses/arxiv-2401.13932-CC-BY-4.0.txt. Characters: every retained excerpt is pure ASCII, so the delivered engine, XeLaTeX with its default Latin Modern text font, reports no missing character.
    \begin{equation}\label{derpt}
        \jdz{ t^k \partial^k_t K(t, x, y) } \leq \frac{c t^{\beta_2} }{(t + |x - y|)^{d + \beta_2}  }.
    \end{equation}

\begin{lemma}\label{BoundTent}
Let $X$ be any fixed Banach space and $1<q<\infty$. Assume that the kernel $\K_{t, s}(x, y)$ satisfies the following estimation: there exist positive constants $ \kappa $, $ \e$, $C$ such that 
\begin{equation}\label{esKts}
    |\K_{t, s}(x, y)| \leq \frac{ C\min\lk{\frac{s}{t},  \frac{t}{s}}^\e \min\lk{ \frac{1}{t}, \frac{1}{s} }^d }{\sk{ 1 + \min\lk{ \frac{1}{t}, \frac{1}{s} }|x - y| }^{d + \kappa}  }.
\end{equation}
    Then the linear operator $\K$ initially defined on $C_c(\real^{d+1}_+)\otimes X$ extends to a bounded linear operator on $T^p_q(\real^{d+1}_+; X)$ for $1 \leq p < \infty$. More precisely, 
    \[
        \| \K(f) \|_{T^p_q(X)} \lesssim_{\e, \kappa} p^{\frac{1}{q}} \| f \|_{T^p_q(X)}, \quad \forall\, f \in T^p_q(\real^{d+1}_+; X), \ 1 \leq p < \infty.
    \]
    Furthermore, for any $f \in C_c(\real^{d+1}_+)\otimes X$, we have
    \[
        \| \mathcal{C}_q\sk{\K(f)} \|_p \lesssim_{\e, \kappa} \| \mathcal{C}_{q}(f) \|_{p}, \quad 1 \leq p \leq \infty.
    \]
\end{lemma}

\begin{lemma}\label{l3.8}
Let $X$ be any fixed Banach space and $1<q<\infty$. The operator $\pi_L$ initially defined on $C_c(\real^{d+1}_+)\otimes X$ extends to a bounded linear operator from $T^p_q(\real^{d+1}_+; X)$ to $H^p_{q, L}(\real^d; X)$ for $1 \leq p < \infty$. More precisely,
    \[
        \| \pi_L(f) \|_{H^p_{q, L}(X)} \lesssim_{\beta} p^{\frac{1}{q}}\| f \|_{T^p_q(X)}, \quad \forall\, f \in T^p_q(\real^{d+1}_+;X), \ 1 \leq p < \infty.
    \]
    Furthermore, for any $f \in C_c(\real^{d+1}_+)\otimes X$, we have
    \[
        \| \pi_L(f) \|_{BMO^p_{q, L}(X)} \lesssim_{\beta} \| \mathcal{C}_q(f) \|_{p}, \quad 1 \leq p \leq \infty.
    \]
\end{lemma}

\begin{proof}
Let $f\in C_c(\real^{d+1}_+)\otimes X$. Recall that $k(t, x, y)$ is the kernel of the operator $\Q$, then
\begin{equation}
\begin{split}
    \Q[\pi_L(f)](x, t) & = \int_{\real^d} k(t, y, z)\pi_L(f)(z) \,\d z  \\
    & = \int_{\real^d} k(t, y, z) \sk{ \int_{\real^{d+1}_+} k(s, z, w)f(w, s) \,\frac{\d w\d s}{s} } \,\d z  \\
    & = \int_{\real^{d+1}_+} \sk{ \int_{\real^d} k(t, y, z) k(s, z, w) \, \d z }f(w, s) \,\frac{\d w\d s}{s}.
\end{split}
\end{equation}
We denote by
\[
    \Phi_{t, s}(y, w) = \int_{\real^d} k(t, y, z) k(s, z, w) \, \d z. 
\]
Note that $k(t, \cdot, \cdot)$ is the kernel of the operator $\Q=-te^{-tL}$, thus $\Phi_{t,s}(\cdot,\cdot)$ is the kernel of $-tLe^{-tL}\circ(-sLe^{-sL})=tsL^2e^{-(t+s)L}$. On the other hand, $\partial^2_r (e^{-rL})|_{r=t+s}=L^2e^{-(t+s)L}$ which has the kernel $\partial^2_r K(r,\cdot,\cdot)|_{r=t+s}$. Then by \eqref{derpt}, we obtain 
\begin{align*}
    |\Phi_{t, s}(y, w)| \lesssim_{d, \beta} \frac{ts}{(t+s)^{2-\beta}(t+s+|y - w|)^{d+ \beta} }  \lesssim_{\beta} \frac{ \min\lk{\frac{s}{t},  \frac{t}{s}} \min\lk{ \frac{1}{t}, \frac{1}{s} }^d }{\sk{ 1 + \min\lk{ \frac{1}{t}, \frac{1}{s} }|x - y| }^{d + \beta}  }.
\end{align*}
Denote by 
\begin{align}\label{Proj}
    \mathcal{P} = 4 \mathcal Q \circ \pi_L. 
\end{align}
 From Lemma \ref{BoundTent}, we conclude that $\mathcal{P}$ initially defined on $C_c(\real^{d+1}_+)\otimes X$ extends to a bounded linear operator on $T^p_q(\real^{d+1}_+; X)$. Moreover,
\[
    \norm{  \mathcal{P}(f) }_{T^p_q(X)} \lesssim_{\beta} p^{\frac{1}{q}}\norm{ f}_{T^p_q(X)}.
\]
Therefore
\[
    \| \pi_L(f) \|_{H^p_{q, L}} = 4^{-1}\| \mathcal{P}(f) \|_{T^p_{q}(X)} \lesssim_{\beta} p^{\frac{1}{q}}\| 
    f \|_{T^p_q(X)},
\]
which is the desired assertion.

For the second part, we obtain the desired assertion from Lemma \ref{BoundTent} immediately.
\end{proof}

\end{document}
